% ---
% id: meat-core/lem:ev-correctness
% kind: lem
% title: "Correctness from error verification"
% summary: >-
%   In the fixed-length post-processing model the two parties accept with
%   different keys with probability at most ε_EV, for every attack and every
%   output length. It is the correctness half of the protocol's security.
% user_statement: >-
%   This is the correctness half of security. Before the key is formed, the
%   key holder announces a ⌈log₂(1/ε_EV)⌉-bit hash of its raw key, computed
%   with a hash function drawn at random independently of the raw keys, and
%   the partner compares it with the hash of its own corrected copy; a
%   mismatch aborts the run. Two different raw keys pass only through a hash
%   collision, so the chance that both parties accept yet hold different keys
%   is at most ε_EV, whatever the attacker does and whatever key length you
%   choose. It needs a two-universal hash family ([[A-HASH]]), seeds
%   independent of the raw keys ([[A-SEED]]) and an implementation that aborts
%   as specified ([[A-CONFORM]]), and takes messages to arrive unaltered;
%   [[cor:composed-security]] adds real authentication. It does not promise
%   that error correction succeeds, only that a failure is caught, except with
%   probability ε_EV.
% prominence: standard
% status: accepted
% uses: [def:postprocessing-model]
% assumes: [A-HASH, A-SEED, A-CONFORM]
% refs:
%   - "tupkary2025survey, Sec. 2.3.1, Def. 2.2, Eq. (4)"
% provenance: newly-derived
% lean:
% ---

The argument reads only the verification step of the model: two different
raw strings pass it only if their short hashes collide, and the hash is
chosen after, and independently of, the strings it compares.

\begin{lemma}\label{lem:ev-correctness}
Let the protocol run the \tref{fixed-length post-processing model}
(Definition~\ref{def:postprocessing-model}) in the
\tref{honest-authentication setting}, let Assumptions~\ref{ass:hash},
\ref{ass:seed} and~\ref{ass:conform} hold, and let $\eev\in(0,1]$. Then for every attack and
every output length $\ell$ the protocol is $\eev$-correct on the accepted
state, in the sense of \cite[Sec.~2.3.1, Def.~2.2, Eq.~(4)]{tupkary2025survey}:
\begin{equation}\label{eq:ev-correct}
  \Pr\big[\Omega_{\mathrm{acc}}\wedge K_A\ne K_B\big]\le\eev .
\end{equation}
\end{lemma}

\begin{proof}
Both parties abort when the verification flag is $\mathrm{rej}$ and output
the keys the model computes (Assumption~\ref{ass:conform}\,(iii)--(iv)). On
$\Omega_{\mathrm{acc}}\wedge\{K_A\ne K_B\}$ the hashed keys differ, so the
raw strings differ, yet their $\lceil\log_2\frac1{\eev}\rceil$-bit
two-universal hashes agreed. (If the raw strings agree, $K_A=K_B$ because
the same public $G$ is applied: in the honest-authentication setting the
partner receives the holder's $G$ unaltered.) The verification function
$h_{\mathrm{EV}}$ is selected by $S_{\mathrm{EV}}$, which
Assumption~\ref{ass:seed} makes independent of the pair
$(Z_1^n,\hat Z_1^n)$, and its family is two-universal
(Assumption~\ref{ass:hash}\,(ii)). Conditioned on any pair of distinct
strings, the probability that their hashes agree is therefore at most
$2^{-\lceil\log_2(1/\eev)\rceil}\le\eev$, and averaging over the pair gives
\eqref{eq:ev-correct}.
\end{proof}
